\documentclass[12pt,a4paper]{report}

% Page geometry
\usepackage{geometry}
\geometry{margin=2.5cm}

% Fonts and Persian support
\usepackage{fontspec}
\usepackage{graphicx}
\usepackage{xcolor}
\usepackage{hyperref}
\usepackage{titlesec} % برای کنترل فاصله‌های عناوین

% Persian support (must be loaded last!)
\usepackage{xepersian}

\hypersetup{
	colorlinks=true,
	linkcolor=blue,
	urlcolor=cyan,
	unicode=true
}

% Font settings
\settextfont[Scale=1.0]{Vazirmatn}
\setlatintextfont{Times New Roman}
\setdigitfont{Times New Roman}

% Layout tweaks
\setlength{\parindent}{0pt}
\linespread{1.3}

% اصلاح فاصله‌ها برای \chapter
\titlespacing*{\chapter}{0pt}{-20pt}{20pt}

\begin{document}
	
	\chapter{مقدمه}
	\vspace{-2.5\baselineskip}
	\noindent\rule{\textwidth}{0.4pt}

	
	  \hspace{-0.1cm} 
	  با گسترش روزافزون اینترنت و افزایش نیاز به خدمات آنلاین، پلتفرم‌های واسطی که ارتباط بین ارائه‌دهندگان خدمات و مشتریان را تسهیل می‌کنند، اهمیت ویژه‌ای پیدا کرده‌اند. در بسیاری از حوزه‌ها، از جمله مدیریت خوابگاه‌های دانشجویی، یافتن یک سیستم مدیریتی قابل اعتماد و کارآمد برای رسیدگی به نیازهای ساکنان می‌تواند چالش‌برانگیز باشد. از سوی دیگر، مدیران خوابگاه‌ها نیز نیاز به یک بستر کارآمد برای مدیریت ثبت‌نام‌ها، نظارت بر درخواست‌های تعمیرات، و برقراری ارتباط مؤثر با دانشجویان دارند.
	  
	  پروژه «سیستم مدیریت خوابگاه» با هدف پاسخ به این نیازها، به عنوان یک پلتفرم آنلاین طراحی و پیاده‌سازی شده است که فرآیند مدیریت خوابگاه، پیگیری درخواست‌ها، و ارتباط بین اعضای مختلف را تسهیل می‌کند. این سیستم به سه دسته اصلی از کاربران خدمت‌رسانی می‌کند: \textbf{دانشجویان} که نیاز به ثبت‌نام در خوابگاه و دسترسی به خدمات مختلف دارند، \textbf{متخصصان سرویس} که در حوزه‌های مختلف تعمیرات و نگهداری تخصص دارند و به دنبال دریافت و انجام درخواست‌های کاری می‌گردند، و \textbf{مدیران} که مسئولیت نظارت، تایید، و کنترل کلیه فرآیندها را بر عهده دارند.
	  
	  هدف اصلی این پروژه، طراحی و پیاده‌سازی یک پلتفرم جامع برای مدیریت هوشمند خوابگاه‌های دانشجویی در یک محیط دیجیتالی امن و کارآمد است. این اهداف به شرح زیر هستند:
	  
	  \begin{itemize}
	  	\item \textbf{تسهیل ارتباط}: ایجاد یک پل ارتباطی مستقیم و کارآمد بین دانشجویان، متخصصان سرویس، و مدیریت خوابگاه.
	  	
	  	\item \textbf{شفافیت و امنیت}: فراهم کردن یک سیستم احراز هویت قوی، تایید ایمیل، و مکانیزم‌های نظارتی برای افزایش اعتماد در پلتفرم.
	  	
	  	\item \textbf{مدیریت هوشمند درخواست‌ها}: ایجاد سیستمی برای ثبت درخواست‌های تعمیرات با امکان بارگذاری تصاویر، تعیین اولویت توسط مدیران، تخصیص به متخصصان مربوطه، و پیگیری وضعیت.
	  	
	  	\item \textbf{مقیاس‌پذیری}: طراحی معماری به گونه‌ای که امکان گسترش سیستم برای افزودن قابلیت‌های بیشتر و پشتیبانی از تعداد بیشتری از کاربران در آینده فراهم باشد.
	  	

	  	\item \textbf{کنترل کیفیت}: پیاده‌سازی سیستم امتیازدهی و بازخورد برای ارزیابی کیفیت خدمات ارائه شده توسط متخصصان سرویس.
	  \end{itemize}


\section{نمودار کلاس سیستم}

\noindent\rule{\textwidth}{0.4pt}

در این بخش، نمودار کلاس سیستم مدیریت خوابگاه ارائه می‌شود که ساختار کلی سیستم و روابط بین اجزای مختلف آن را نشان می‌دهد. این نمودار شامل کلاس‌های اصلی سیستم از جمله مدیریت کاربران، درخواست‌های تعمیرات، سیستم تیکت‌ها و ساختار فیزیکی خوابگاه است.

\begin{figure}[h!]
	\centering
	\resizebox{1\textwidth}{!}{%
		\includegraphics{../issue/classes.png}%
	}
	\caption{\lr{classes of dormitory system}}
	\label{fig:websocket}
\end{figure}

همان‌طور که در شکل  مشاهده می‌شود، سیستم حول کلاس \texttt{User} به عنوان کلاس مرکزی طراحی شده است که انواع مختلف کاربران را پشتیبانی می‌کند. این کلاس دارای دو کلاس فرعی اصلی است:

\begin{itemize}
	\item \textbf{Student}: برای مدیریت دانشجویان ساکن در خوابگاه
	\item \textbf{ServiceExpert}: برای مدیریت متخصصان ارائه‌دهنده خدمات تعمیرات
	\item \textbf{Admin}: برای ادمین های نظاره کننده بر سیستم
\end{itemize}

\hspace{0.1cm}کلاس‌های \texttt{Building} و \texttt{Room} ساختار فیزیکی خوابگاه را مدل‌سازی کرده‌اند، در حالی که کلاس \texttt{MaintenanceRequest} فرآیند درخواست و انجام تعمیرات را مدیریت می‌کند. سیستم تیکت‌ها نیز از طریق کلاس‌های \texttt{Ticket} و \texttt{TicketMessage} برای برقراری ارتباط در زمان واقعی بین کاربران پیاده‌سازی شده است.


\end{document}
